\section{Applicazioni gestionali e DBMS}
Le \textbf{Applicazioni} informatiche in ambito gestionale sono nate negli anni '70 e sono caratterizzate da algoritmi di elaborazione semplici e una
grande quantità di dati condivisi da più applicazioni, ambo derivanti dall'ambiente in cui sono introdotte, chiamato \textbf{sistema informativo}.
Quest'ultimo è definito formalmente come l'\textit{insieme delle attività umane e dei dispositivi di memorizzazione ed elaborazione atte alla gestione e
all'organizzazione di informazioni di interesse per un'organizzazione}. Il sistema informativo è composto a sua volta da due componenti: \begin{itemize}
    \item \textbf{Dato}:  Elemento di conoscenza di base, costituito da simboli da elaborare.
    \item \textbf{Informazione}: Interpretazione dei dati che consente di ottenere conoscenza di fatti e situazioni.
\end{itemize}

\noindent È possibile studiare il comportamento dei sistemi informativi mediante diagrammi di flusso, i quali \textbf{definiscono} archivi, sorgenti,
utenti, procedure e flussi di dati, quindi rappresentano le informazioni come collezioni di dati. Questi elementi portano a quella che è chiamata
\textbf{soluzione convenzionale}, che mostra i vari programmi interagire col file system, il quale mantiene i dati. Tuttavia presenta alcuni problemi,
come per esempio un accesso lento ai dati e la loro ridondanza.\par
Un'altra soluzione pone come middleware posto tra le applicazioni e il file system un \textbf{Database Management System} (DBMS), un sistema che
gestisce le \textbf{basi di dati} su memoria secondaria, assicurando affidabilità, privatezza e accesso efficiente. Presenta inoltre due vantaggi
principali: \begin{itemize}
    \item \textbf{Maggiore astrazione}: I dati sono visti come oggetti o righe di tabelle quando il DBMS si interfaccia con le applicazioni.
    \item \textbf{Maggiore efficacia}: L'accesso ai dati è basato su un linguaggio di interazione.
\end{itemize}

\noindent In particolare, i \textbf{linguaggi di interazione} si dividono in due categorie: la prima è atta alla \textbf{definizione dei dati} (DDL),
mentre la seconda concerne la \textbf{manipolazione dei dati} (DML), a sua volta suddivisa in \textbf{linguaggi di interrogazione}, per estrarre
informazioni dal DB,  e \textbf{di manipolazione}, per popolare il DB o modificarne i dati. Il DBMS fornisce inoltre un \textbf{Modello dei Dati},
l'insieme dei costrutti per la descrizione di struttura e proprietà dell'informazione all'interno del DB. Ad oggi, vengono utilizzati i seguenti
modelli: \begin{itemize}
    \item Modello relazionale: SQL92 o SQL2
    \item Modello a oggetti
    \item Modello object-relational: SQL99 o SQL3
    \item Modello basato su documenti: JSON
    \item Modelli NoSQL: Basati su colonne, costrutti chiave-valore
\end{itemize}

\noindent Bisogna tuttavia fare attenzione a distinguere il modello dei dati dallo \textbf{schema del DB}, la descrizione di struttura e proprietà della
base di dati che fa uso dei costrutti del modello, e dall'\textbf{istanza del DB}, costituita dagli effettivi valori che popolano le strutture dati in
una determinata istanza. In soldoni, il modello definisce i costrutti, lo schema la loro relativa implementazione e l'istanza una situazione specifica
con determinati valori.\par
Passiamo ora all'architettura generale dei DBMS: presentano tre livelli: \textbf{esterno}, che descrive una porzione dello schema logico per
un'interfaccia specifica, \textbf{logico}, la rappresentazione di struttura e proprietà del DB, definite con il modello dei dati, e \textbf{interno},
rappresentante il DB mediante le strutture fisiche di memorizzazione. In un contesto tale è necessario chiarire che le basi di dati devono mantenere due
tipi di indipendenza: \begin{itemize}
    \item \textbf{Fisica}: Schema logico è indipendente da quello interno, facendo sì che la variazioni sul secondo non impattino il primo.
    \item \textbf{Logica}: Schemi esterni indipendenti da quello logico, facendo sì che le variazioni nel secondo, purché non rimuovano dati, non
	impattano i primi.
\end{itemize}

%

\section{Creazione di una base di dati}
Per progettare adeguatamente una base di dati è necessario studiare il ciclo di vita del processo di automazione all'interno di un sistema informativo.
In primo luogo è richiesto uno \textbf{studio di fattibilità} per definire costi ed eventuali alternative; segue la \textbf{raccolta e analisi dei 
requisiti} del sistema, producendo una descrizione completa ma informale, per poi terminare con l'effettiva \textbf{progettazione}. Quest'ultima fase è
suddivisa in due attività: \begin{itemize}
    \item \textbf{Progettazione dei dati}: Descrizione formale dei dati attraverso lo schema, comprende l'implementazione del DBMS seguita dalla
	validazione e collaudo.
    \item \textbf{Progettazione delle applicazioni}: Descrizione formale delle applicazioni attraverso le specifiche.
\end{itemize}

\noindent In questo corso ci concentreremo sulle metodologie di progettazione dei dati. Un buon metodo è costituito da una decomposizione iniziale in
passi dell'attività di progetto, attuazione di strategie e modelli di riferimento da seguire, come anche di criteri di scelta da applicare. La
metodologia deve inoltre essere generale, semplice da usare, ma soprattutto in grado di produrre un risultato completo e corretto.\par
Nel concreto, la progettazione dei dati presenta tre fasi principali: \textbf{concettuale}, per rappresentare il contenuto informativo del DB
formalmente e indipendente da tecnologia e operazioni, portando alla creazione dello schema concettuale, \textbf{logica} per tradurre lo schema
concettuale in quello logico, aderente al modello dei dati del DBMS scelto, e \textbf{fisica}, per completare lo schema logico coi parametri relativi
alla memorizzazione fisica dei dati e metodi di accesso per garantire efficienza, fornendo lo schema fisico.

%

\section{Progettazione concettuale}
Il modello storicamente usato per la progettazione concettuale delle basi di dati è quello \textbf{Entità-Relazione} (ER); propone un insieme di
strumenti grafici formali e indipendenti dalla tecnologia\footnote{Privi di qualsivoglia ottimizzazioni o semplificazioni che possono avere luogo in
base agli strumenti scelti.}, detti \textbf{costrutti}, per specificare struttura e proprietà dei dati da rappresentare al suo interno.\par
La definizione di ogni costrutto si ha specificando il loro significato, sintassi grafica e relative istanze; quelli di base sono: \begin{itemize}
    \item \textbf{Entità}: Classe di oggetti con attributi comuni, esistenza autonoma e identificazione univoca, rappresentate con un rettangolo
    contenente il loro nome.
    \item \textbf{Relazione}: Legame logico tra due o più entità, rappresentate da un rombo a cui si collegano attraverso linee le entità coinvolte.
    Il nome è posto vicino alla forma geometrica.
    \item \textbf{Attributo}: Proprietà elementare di un'entità o una relazione; associa a ogni loro istanza un singolo valore appartenente a un
    dominio, indicandolo con una piccola riga alla cui estremità è posto un pallino vuoto con a fianco il suo nome.
\end{itemize}

\noindent Ognuno di questi costrutti si può \textbf{istanziare}, ovvero concretizzare, all'interno della base di dati; sono annotati con
$I(nomeCostrutto)$ e i relativi insiemi sono detti \textbf{popolazioni}. Alla nascita della base di dati ogni popolazione è vuota, poiché nessun
costrutto è stato ancora istanziato. Istanziare un'entità significa aggiungere un oggetto alla base di dati e farlo per un attributo corrisponde ad
associare un valore del dominio che assume a un'istanza di entità o relazione. Per queste ultime in particolare, definiamo le loro istanze come
n-uple di istanze di entità; formalmente: \[
    (e_1, ..., e_n)\text{, con } e_i\in I(E_i), 1 \leq i \leq n
\]

\noindent E la relativa popolazione rappresenta l'insieme delle coppie delle due entità coinvolte relazionate in un certo istante di tempo. Si tratta
di una rappresentazione univoca, quindi non è possibile annotare più n-uple allo stesso modo.

\begin{comment}
% TODO Inserisci immagini entità, relazione

Lez5 Basi di dati - 15-10-26

- Esercizio ER1, ricava il drawio

--- Adesso ha iniziato le slide ER2
Necessario identificare quali entità sono corrette e quali non. Nelle relazione non puoi usare questo costrutto seguente.
- Identificatore di entità: Data un'entità E, è un insieme di proprietà (attributi e/o relazioni) che identificano univocamente le istanze dell'entità.
	Le proprietà identificano univocamente le istanze di E se non esistono due istanze di E che presentano gli stessi valoru/istanze nelle proprietà dell'insieme.
	Si rappresenta graficamente in più modi:
		- Se un identificatore è costituito da un solo attributo, si riempie il pallino di colore.
		- (Identificatori interni) Se è costituito da più attributi va tirata una riga tra i due con all'estremità il pallino pieno. Attenzione che senza la riga, si sta dicendo che
			due identificatori stanno identificando qualcosa di diverso.
		- (Identificatori esterni) Se include una relazione, quest'ultima deve comportarsi come una funzione, quindi essere una singola istanza per identificare l'oggetto.
	Gli identificatori implicano anche un tipo di dipendenza dall'identificatore, perché senza di essi non si avrebbe modo di capire cosa sono. Non è per nessun motivo possibile
	dare un identificatore a un'istanza di relazione.

\end{comment}

%

\section{Progettazione logica}